% Generated draft: requires mathematical and source review.
\documentclass[11pt]{article}
\usepackage[margin=25mm]{geometry}
\usepackage{fontspec}
\setmainfont{Latin Modern Roman}
\usepackage{amsmath,amssymb,mathtools,unicode-math}
\setmathfont{Latin Modern Math}
\usepackage{enumitem,needspace,xurl,hyperref}
\setlength{\parindent}{0pt}
\setlength{\parskip}{5pt}
\setlength{\emergencystretch}{3em}
\newcommand{\N}{\mathbb{N}}
\newcommand{\Z}{\mathbb{Z}}
\newcommand{\Q}{\mathbb{Q}}
\newcommand{\R}{\mathbb{R}}
\newcommand{\C}{\mathbb{C}}
\newcommand{\F}{\mathbb{F}}
\newcommand{\A}{\mathbb{A}}
\newcommand{\PP}{\mathbb{P}}
\newcommand{\E}{\mathbb{E}}
\newcommand{\eps}{\varepsilon}
\newcommand{\dd}{\,\mathrm{d}}
\begin{document}
\section*{UnsolvedMath ID: 30005973}
\textbf{Source problem code:} catalogue entry
\subsection{Definitions and mathematical statement}
A smooth plane quartic over a field \(K\) is a smooth geometrically connected curve \(C\subset\PP^2_K\) defined by a homogeneous polynomial of degree \(4\); it has genus \(3\). In the usual complete-intersection description, a Cayley octad consists of eight distinct geometric points in \(\PP^3\) forming the scheme-theoretic intersection of three independent quadrics. The supplied record refers to a Cayley octad associated with \(C\), but does not specify the construction, its auxiliary choices, or its hypotheses.

In a discrete valuation setting, a stable model after a finite extension is a proper flat family with generic fibre the extended curve and geometric special fibre a connected nodal curve satisfying \(2g_v-2+n_v>0\) on every irreducible component. Here \(g_v\) is the genus of the normalization of that component and \(n_v\) counts branches of nodes on it. One possible meaning of stable reduction type is the isomorphism class of the dual graph of this special fibre, with vertices weighted by \(g_v\), edges representing nodes, and loops allowed. The record does not specify whether this is the intended meaning or whether additional data are retained.

\textbf{Conjecture as supplied, with unresolved scope.} The combinatorial configuration data of an associated Cayley octad determine the stable reduction type of a smooth plane quartic. More formally, let \(\mathscr D\) denote the intended class of admissible curve--octad data, let \(\mathcal P(C,O)\) denote the intended combinatorial invariant, and let \(T(C)\) denote the intended stable reduction type. The asserted determination property has the form
\[
 \forall (C_1,O_1),(C_2,O_2)\in\mathscr D,\qquad
 \mathcal P(C_1,O_1)\cong\mathcal P(C_2,O_2)
 \ \Longrightarrow\ T(C_1)\cong T(C_2).
\]
This is a formulation template, not a fully specified conjecture: the source does not define \(\mathscr D\), \(\mathcal P\), the relevant equivalences, or \(T\). In particular, it leaves the valued base field, residue characteristic restrictions, valuation-dependent configuration data, and treatment of choices of associated octad unresolved. No assertion over all fields or residue characteristics is made here.

The source motivates the conjecture by replacing the eight Weierstra\ss{} points in the hyperelliptic comparison with a Cayley octad, analogously to the role of a hyperelliptic cluster picture. The hypotheses of that comparison are also unspecified.
\subsection{Short English statement}
The proposed conjecture says that suitable combinatorial data attached to the eight points of an associated Cayley octad determine a smooth plane quartic's stable reduction type. The supplied record does not define those data or the precise arithmetic scope, so a complete mathematical formulation requires further source clarification.
\subsection{Sources}
\begin{itemize}
\item Supplied local UnsolvedMath registry record, permanent ID \(30005973\), title \emph{Stable Reduction of Plane Quartics from Cayley Octads}, source problem code \texttt{catalogue entry}, in the Oberwolfach Reports Open Problems collection. No individual proposer is supplied; the proposed year is recorded as \(2024\).
\item \emph{Explicit Methods in Number Theory (2025)}, Oberwolfach Reports, pp.\ 2309--2368, DOI \texttt{10.4171/owr/2024/40}: \url{https://doi.org/10.4171/owr/2024/40}. This is the supplied report provenance and is also listed among the supplied literature sources. The record describes its original statement as a verbatim contextual passage.
\item van Bommel--Docking--Dokchitser--Lercier et al., \emph{Reduction of Plane Quartics and Cayley Octads}: \url{https://doi.org/10.1007/s10208-025-09704-y}. This is supplied as the primary-paper reference; no conjecture number or exact supporting passage is provided.
\item Literature assessment dated exactly \texttt{2026-08-22}: the supplied curation labels the problem open and reports extensive computational and theoretical evidence for the octad-picture criterion, while stating that the 2025 primary paper still presents the stable-reduction correspondence as a conjecture. This is supplied dated curation, not an independently verified literature assessment; its time-sensitive status and cited claims require independent checking.
\item The registry labels the statement \texttt{corrected\_verified}, with the qualification that only formula delimiters, notation, grammar, and typography were repaired and that mathematical content and scope were preserved. This supplied editorial label does not establish mathematical correctness or resolve missing hypotheses.
\end{itemize}
\end{document}
